\documentclass[aspectratio=169,11pt]{beamer}
\usetheme[progressbar=frametitle,numbering=fraction,block=fill]{metropolis}
\usepackage{booktabs}
\usepackage{array}
\newcolumntype{L}[1]{>{\raggedright\arraybackslash}p{#1}}
\usepackage{amsmath,amssymb}
\usepackage{tikz}
\usetikzlibrary{arrows.meta,positioning,calc,shapes.geometric,decorations.pathreplacing}
\usepackage{pgfplots}
\pgfplotsset{compat=1.17}
\usepackage{tcolorbox}

% Palette: deep sea navy, teal, amber accent
\definecolor{sea}{HTML}{0B2545}
\definecolor{teal}{HTML}{13838F}
\definecolor{foam}{HTML}{E6F2F3}
\definecolor{amber}{HTML}{E09F3E}
\definecolor{mute}{HTML}{8A9BA8}
\setbeamercolor{normal text}{fg=sea!90!black,bg=white}
\setbeamercolor{alerted text}{fg=amber!85!black}
\setbeamercolor{frametitle}{bg=sea,fg=white}
\setbeamercolor{progress bar}{fg=teal,bg=foam}
\setbeamercolor{title separator}{fg=teal}
\setbeamercolor{background canvas}{bg=white}
\setbeamercolor{block title}{bg=foam,fg=sea}
\setbeamercolor{block body}{bg=foam!45}
\setbeamercolor{section in toc}{fg=sea}
\makeatletter\setlength{\metropolis@frametitle@padding}{1.6ex}\makeatother

\newcommand{\gapbox}[3]{%
  \begin{tcolorbox}[colback=foam!60,colframe=foam!60,arc=2mm,boxrule=0pt,left=2mm,right=2mm,top=0.6mm,bottom=0.6mm,before skip=0.8mm,after skip=0.8mm]
  {\bfseries\color{teal}#1}\hspace{1mm}{\small\bfseries #2}\\
  {\footnotesize #3}
  \end{tcolorbox}}
\newcommand{\stat}[2]{%
  \begin{minipage}[t]{\linewidth}\centering
  {\fontsize{22}{24}\selectfont\bfseries\color{teal}#1}\\[1mm]
  {\small #2}
  \end{minipage}}

% Evidence labels: what each statement rests on
\newcommand{\evtag}[2]{\tikz[baseline=(t.base)]\node[fill=#1,text=white,rounded corners=0.7mm,inner xsep=1mm,inner ysep=0.35mm,font=\tiny\bfseries](t){#2};}
\newcommand{\claim}{\evtag{sea!75}{CLAIM}}
\newcommand{\sugg}{\evtag{teal}{SUGGESTED}}
\newcommand{\open}{\evtag{amber}{OPEN}}
\newcommand{\src}{\evtag{mute}{SOURCE}}

% Symbols used throughout
\newcommand{\Nsym}{N_{\mathrm{sym}}}
\newcommand{\Kup}{K_{\mathrm{upd}}}
\newcommand{\dtol}{d_{\mathrm{tol}}}
\newcommand{\etol}{e_{\mathrm{tol}}}
\newcommand{\Pset}{\mathcal{P}}

\title{Doppler-Aided Link Adaptation\\for Underwater Acoustic Links}
\subtitle{Main idea, gaps, methodology and theoretical objectives --- for group discussion}
\date{8 October 2026 (revised 9 October 2026 after the prior-art search)}
\institute{Target: IEEE JSAC special issue on Next-Generation Underwater Acoustic Communication Systems (deadline 31 March 2027)\\[1mm]
\claim\ stated in the idea document, not yet established \quad \sugg\ design option added after review \quad \open\ needs material, a proof or an experiment \quad \src\ found by targeted search, not yet checked}

\begin{document}

\maketitle

% ===========================================================
\section{Main idea}

\begin{frame}{The proposal in one slide}
\begin{columns}[T]
\begin{column}{0.61\textwidth}
\small
Every packet, data or ACK, carries a free measurement of the \alert{Doppler scale} $\epsilon=v_{\mathrm{rel}}/c$, the same at both ends. Both ends fuse them (and motion sensors, if any) into one \alert{kinematic forecast} of the drift $d=\dot\epsilon$ and its uncertainty, which drives four blocks:
{\scriptsize
\begin{itemize}\setlength{\itemsep}{0pt}\setlength{\topsep}{2pt}
\item \textbf{frame}: extra pilots where the forecast is uncertain, over a coherence-limited grid; packet length, modulation
\item \textbf{receiver}: pre-warped along the forecast, tracking only its error
\item \textbf{transmitter}: bulk Doppler pre-compensated, so several nodes can share the uplink at once
\item \textbf{decision}: threshold rule on the predicted PER, margin from delayed ACKs; uplink guard bands and admission from $\sigma_e$
\end{itemize}}

{\scriptsize \textbf{Candidate claim} \sugg: Only a forecast, learned from the link's own two-way Doppler, lets a transmitter pre-compensate, so several nodes share an uplink (NTN: from ephemeris). Bandit: follow-up.}
\end{column}
\begin{column}{0.36\textwidth}
\begin{tcolorbox}[colback=sea,colframe=sea,arc=2mm,coltext=white,before skip=0pt]
\footnotesize\raggedright
\textbf{Scope}\\[1mm]
Delay-tolerant traffic\\
Mobile or surface-coupled nodes\\
Seas with a swell component\\
Paths resolvable in delay or Doppler\\
One link, or $K$ nodes to a gateway ($K=1$: one link)\\[1mm]
\textit{Still, fixed links: fall back to the baseline; ``no harm'' once overheads are counted is still to be shown.}
\end{tcolorbox}
\end{column}
\end{columns}
\end{frame}

% ===========================================================
\section{Gaps with the literature}

\begin{frame}{Related work (1): adapting the link despite stale feedback}
\scriptsize\setlength{\tabcolsep}{3pt}\renewcommand{\arraystretch}{0.92}
\begin{tabular}{@{}L{0.21\textwidth}L{0.46\textwidth}L{0.29\textwidth}@{}}
\toprule
\textbf{Line of work} & \textbf{Representative papers} & \textbf{Difference}\\
\midrule
Physics-informed, uncertainty-aware AMC (closest) & Anjangi et al., UComms 2018 (Doppler spread as AMC feature); Wu, Chitre, Anjangi, JOE 2025 (GP interval BER + feedback scheduling) & MCS only; no kinematic forecast, packet layout or timing\\
Channel prediction for AMC & Radosevic et al., OCEANS 2011, JOE 2014; Tian et al., OCEANS 2026; Wang, OCEANS 2022; Jing, IoT-J 2025 (delay-Doppler) & Predict the channel, not the motion; MCS or bit loading\\
Prediction instead of feedback & Park et al., Remote Sens.\ 2026 (closest to our decision rule) & Doppler index, margin, PER target; Raptor overhead only\\
Feedback timing & AoI-aware hierarchical bandit (Busacca et al., GLOBECOM 2025) & No physical model\\
Error control under delay & OLLA (LTE/5G); ACI (Gibbs--Cand\`es 2021); delayed ACI (El Halabi--Brandt 2026) & \textcolor{amber!85!black}{Same update rule as our margin}; new: action-dependent delay, service accounting\\
AMC survey and benchmark & Cooper-Baldock et al., Appl.\ Ocean Res.\ 2026 (open-data benchmark) & Baselines to compare against\\
\bottomrule
\end{tabular}

\vspace{0.5mm}
{\color{mute} Details: \texttt{literature/prior\_art\_summary.md}.}
\end{frame}

\begin{frame}{Related work (2): Doppler is compensated, and packets are designed offline}
\scriptsize\setlength{\tabcolsep}{3pt}\renewcommand{\arraystretch}{0.9}
\begin{tabular}{@{}L{0.21\textwidth}L{0.45\textwidth}L{0.30\textwidth}@{}}
\toprule
\textbf{Line of work} & \textbf{Representative papers} & \textbf{Use of Doppler}\\
\midrule
Physics of wave-induced Doppler & Walstead, JASA 2016--17; Rawat, IEEE TAES 2021 (Doppler spectrum from sea-surface motion); Li et al., JOE 2008 (residual offset varies with waves) & Characterisation; design of compensation\\
Packet structure for a motion range & Zakharov et al., JOE 2019 ($|a|\le0.25$\,m/s$^2$ design, $\le1$ simulated; $\approx$1\,s packets) & \textcolor{amber!85!black}{Fixed offline design}, not adapted online\\
Acceleration-aware compensation & Abdelkareem et al., Ad Hoc Netw.\ 2016 (previous packet gates the search); Cho--Kim, IET Commun.\ 2026; Hu et al., MSN 2024 (Kalman) & Absorbs constant drift; models stay inside one packet\\
Doppler prediction in the receiver & Ceballos Carrascosa and Stojanovic, JOE 2010 (Doppler predicted across OFDM blocks) & In the receiver; \textcolor{amber!85!black}{not used by the transmitter}\\
Doppler of both ends & DE-Sync 2018, APE-Sync 2019 (returned in the reply); Aubry 2019, Choi 2022 (navigation) & Clock sync, position; not the link\\
\bottomrule
\end{tabular}

\vspace{0.5mm}
\scriptsize\textbf{Enablers now public:} Underwater Acoustic Channel Library (Li et al., 2026; bulk drift kept, one-way, $\approx$50--60\,s records), Qarabaqi--Stojanovic simulator, NOAA NDBC wave spectra. \alert{Caution:} Watermark may remove bulk Doppler.
\end{frame}

\begin{frame}{Related work (3): the closest analogues are outside underwater acoustics}
\scriptsize\setlength{\tabcolsep}{3pt}\renewcommand{\arraystretch}{0.92}
\begin{tabular}{@{}L{0.30\textwidth}L{0.34\textwidth}L{0.32\textwidth}@{}}
\toprule
\textbf{Work} & \textbf{Shared mechanism} & \textbf{What it lacks}\\
\midrule
Yuan et al., JSTSP 2021 (V2I ISAC) & Prediction drives TX, pilots and RX: \alert{the concept} & Constant velocity; no two-way Doppler, PER layer\\
Simeone--Spagnolini, ICC 2004 & Pilots from Kalman covariance & Fading, not kinematics\\
Li et al., arXiv 2606.16717 & Wave-driven buoy state model & RF beamforming\\
3GPP NTN (TS 38.300) & Continuous Doppler pre-compensation & External ephemeris\\
DE-Sync 2018, APE-Sync 2019 & Far end's Doppler in the reply & Clock skew only\\
\bottomrule
\end{tabular}

\vspace{1mm}
\footnotesize The architecture is not new in general: claim the underwater instance and its analysis (R1--R4).\\ \open\ Scholar cited-by of Yuan 2021, DE-Sync, APE-Sync not yet checked.
\end{frame}

\begin{frame}{Related work (4): multi-user Doppler is handled at the receiver}
\scriptsize\setlength{\tabcolsep}{3pt}\renewcommand{\arraystretch}{0.92}
\begin{tabular}{@{}L{0.33\textwidth}L{0.33\textwidth}L{0.30\textwidth}@{}}
\toprule
\textbf{Work} & \textbf{Mechanism} & \textbf{What it lacks}\\
\midrule
Daoud et al., IEEE TVT 2015 & Multiple resampling at the receiver, one branch per user & No prior; receiver only\\
Tu, Duman, Stojanovic, Proakis 2011, 2013 & OFDMA; one resampling branch per user at the receiver & $K$ branches; interference remains\\
Zhang et al., arXiv 2026 & Concurrent random access; per-user power control & Path-specific Doppler left in the receiver\\
Hasan et al., Ad Hoc Netw.\ 2025 & Guard bands adapted to mobility (metaheuristic) & No forecast, no pre-compensation\\
Oiwa et al., IEICE Tech.\ Rep.\ 2020 & Accelerometer-aided Doppler compensation & Own node only; not peer reviewed\\
3GPP NTN (TS 38.300) & UE pre-compensates Doppler before sending & Ephemeris, not learned from the link\\
\bottomrule
\end{tabular}

\vspace{1mm}
\footnotesize Library, web, Scholar and 159 citing works checked (9 Oct 2026): no UWA transmitter pre-compensates its own Doppler. Say ``transmit-side'': UWA papers use ``Doppler pre-compensation'' for a receiver step.
\end{frame}

\begin{frame}{Candidate gaps, re-ranked after the library check}
\gapbox{G1}{No packet places training from the error variance of a kinematic Doppler forecast}{Pilots follow mobility \emph{statistics} (NR DMRS, 802.11bd) or spreads (Mangione 2021); Kalman-covariance pilots: radio only (Simeone--Spagnolini 2004).}
\gapbox{G2}{No receiver carries a Doppler forecast with uncertainty across packets}{Kinematic models stay inside one packet (Cho--Kim 2026, Hu 2024); the previous packet only gates the search (Abdelkareem 2016).}
\gapbox{G3}{Doppler from both ends is not tracked as a kinematic process for the link}{Exchanged for clock sync (DE-Sync, APE-Sync); fed one way to navigation (Aubry 2019, Choi 2022).}
\gapbox{G5}{No UWA transmitter pre-compensates Doppler for a multi-user uplink}{Per-user Doppler is separated at the receiver, one resampling branch per user (Tu et al.\ 2011, 2013); guard bands follow mobility (Hasan 2025).}
{\tiny\color{mute} G4 (learning with arm-dependent delay): follow-up paper. Dropped: send timing to the swell.}
\end{frame}

% ===========================================================
\section{Proposed methodology}

\begin{frame}{System at a glance}
\centering
\begin{tikzpicture}[
  box/.style={fill=foam,rounded corners=2mm,minimum height=10mm,text width=2.9cm,align=center,font=\small},
  key/.style={fill=teal,text=white,rounded corners=2mm,minimum height=12mm,text width=3.0cm,align=center,font=\small},
  arr/.style={-{Stealth[length=2.5mm]},thick,sea},
]
\node[box] (tx) at (0,0) {\textbf{Node A}\\ \scriptsize measures $\hat\epsilon$ on ACKs};
\node[box] (rx) at (9.2,0) {\textbf{Node B}\\ \scriptsize pre-warp (1); $\hat\epsilon$ on data};
\draw[arr] ([yshift=2.5mm]tx.east) -- node[above,font=\scriptsize]{data: $a=(M,\Nsym,\Pset)$ + forecast} ([yshift=2.5mm]rx.west);
\draw[arr,teal] ([yshift=-2.5mm]rx.west) -- node[below,font=\scriptsize,align=center]{ACK: outcome of an earlier packet\\B's time-stamped $\hat\epsilon$ ($\approx$20 bits)} ([yshift=-2.5mm]tx.east);
\node[key] (kf) at (0,-2.9) {Kalman filter (3)\\ \scriptsize EKF; Doppler from both ends};
\node[key] (cfg) at (4.6,-2.9) {Pilot placement (2)\\ + threshold rule};
\node[key] (mg) at (9.2,-2.9) {Safety layer\\ \scriptsize PER margin $m_t$ from ACKs};
\draw[arr] (tx.south) -- node[left,font=\scriptsize,align=right]{both\\series} (kf.north);
\draw[arr] (kf) -- node[above,font=\scriptsize,align=center]{forecast\\$\hat d$, $\sigma_e$} (cfg);
\draw[arr] (mg) -- node[above,font=\scriptsize]{$m_t$} (cfg);
\draw[arr] (rx.south) ++(0,-0.15) -- node[right,font=\scriptsize]{outcomes} (mg.north);
\draw[arr,amber] (cfg.north) -- ++(0,0.75) -| ([xshift=6mm]tx.south);
\end{tikzpicture}

\medskip
\footnotesize\raggedright
Shown for one link ($K=1$); numbers in brackets: blocks on the following slides; the uplink with $K$ nodes adds Block 4. \open\ For the technical specification: header fields and quantisation; measurement ordering (out of sequence only with pipelined packets); SNR source; NACK or timeout, packet matching, ACK vs data loss; whether $\hat\epsilon$ needs payload decoding or only preamble acquisition.
\end{frame}

\begin{frame}{Two quantities that must not be confused}
\begin{columns}[T]
\begin{column}{0.5\textwidth}
\small\setlength{\abovedisplayskip}{3pt}\setlength{\belowdisplayskip}{3pt}
\textbf{Doppler scale} (dimensionless)
\[
\epsilon(t)=\frac{v_{\mathrm{rel}}(t)}{c}
\]
\textbf{Physical Doppler drift} (units s$^{-1}$)
\[
d(t)=\dot\epsilon(t)=\frac{\dot v_{\mathrm{rel}}(t)}{c}
\]
{\footnotesize Some papers call $\epsilon$ the ``Doppler rate''. Here $d$ is always the \emph{physical} drift, never the receiver's residual.}
\end{column}
\begin{column}{0.46\textwidth}
\small
\begin{itemize}
\item What harms a packet is what \alert{remains after compensation}: it depends on receiver and action, not on $d$ alone:\\
{\footnotesize forecast error $e$ $\to$ receiver, action $a$ $\to$ residual impairment $\phi_a(e)$ $\to$ PER}
\item Example: $|d|=10^{-4}\,$s$^{-1}$ is a relative acceleration of 0.15\,m/s$^2$ at $c=1500$\,m/s.
\item One \alert{common (bulk)} component is an initial restriction. If path powers, delay spread, differential Doppler or tracking state change PER at equal $d$, that is a model limitation \open.
\end{itemize}
\end{column}
\end{columns}
\end{frame}

\begin{frame}{Model}
\footnotesize\setlength{\abovedisplayskip}{1pt}\setlength{\belowdisplayskip}{1pt}
\textbf{Motion} \claim: $v_{\mathrm{rel}}(t)=$ slow platform drift $+$ damped oscillators (swell, wind sea) learned online.

\vspace{0.6mm}
\textbf{Physical drift} $d=\dot v_{\mathrm{rel}}/c$ in s$^{-1}$, common (bulk) component only: an initial restriction \open.

\vspace{0.6mm}
\textbf{Observations:} time-stamped $\hat\epsilon$ on the preamble of every packet received at \emph{either} end (accuracy set by preamble length and bandwidth); ACK or timeout for earlier packets.

\vspace{0.6mm}
\textbf{Forecast error} $e(u)=d(u)-\hat d(u)$ over the receive interval, standard deviation $\sigma_e(u)$ from the estimator. A receiver without the prior has $\hat d=0$, so $e=d$ (the baseline).

\vspace{0.6mm}
\textbf{Failure} \claim: the receiver turns $e$ into a residual impairment $\phi_a(e)\ge0$, depending on $(M,\Nsym,\Pset)$; a monotone scalar curve is a hypothesis, checked by calibrating each receiver in simulation:
\[
\mathrm{SNR}_{\mathrm{eff}}=\frac{\mathrm{SNR}}{1+\mathrm{SNR}\,\phi_a(e)}<\gamma_a
\iff
\phi_a(e)>\frac{1}{\gamma_a}-\frac{1}{\mathrm{SNR}},
\qquad
\etol(a,\mathrm{SNR})=\sup\Big\{|e|:\phi_a(|e|)\le\tfrac{1}{\gamma_a}-\tfrac{1}{\mathrm{SNR}}\Big\}
\]
Empty set (e.g.\ $\mathrm{SNR}<\gamma_a$): $a$ excluded; unbounded: $\etol=\infty$; non-monotone $\phi_a$: use the set itself. $\gamma_a$: SNR at which $a$ meets the PER target, a measured threshold \sugg.

\vspace{0.6mm}
\textbf{Traffic:} bursty arrivals; waiting $h\le h_{\max}$ set by the MAC; a delay budget per packet.
\end{frame}

\begin{frame}{A minimal action space \sugg}
\begin{columns}[T]
\begin{column}{0.5\textwidth}
\small
\[
a=(M,\ \Nsym,\ \Pset)
\]
\begin{itemize}
\item $M$: modulation order, QPSK or 16-QAM
\item $\Nsym$: OFDM symbols in the packet
\item $\Pset\subset\{1,\dots,\Nsym\}$: positions of full training symbols, $N_{\mathrm{tr}}=|\Pset|$
\end{itemize}
\footnotesize
The receiver pre-warp is not an action: it always follows the forecast. Fixed in the first paper: useful symbol duration $T_s$, cyclic prefix $T_{cp}$, code. Waiting time $h$: set by the MAC.
\end{column}
\begin{column}{0.46\textwidth}
\small
\textbf{Packet duration}
\[
T_a=T_{\mathrm{pre}}+\Nsym\,(T_s+T_{cp})
\]
\textbf{Payload}
\[
B_a=\big(\Nsym-|\Pset|\big)\,N_{\mathrm{sc}}\,\log_2 M\,R_c
\]
{\footnotesize Uniform $\Pset$ with spacing $\Kup$ is the baseline. Exact count from the frame (first estimate from the preamble? final training symbol?) \open. $B_a$: information bits before header, CRC, termination; split messages count every segment and its delay.}
\end{column}
\end{columns}
\end{frame}

\begin{frame}{When the packet is received, and what is known then}
\centering
\begin{tikzpicture}[x=1.3cm,y=1cm,
  lab/.style={font=\scriptsize,align=center},
  arr/.style={-{Stealth[length=2mm]},thick}]
\draw[mute,thick] (-0.3,1.6) -- (9.6,1.6) node[right,font=\scriptsize,text=sea]{Tx};
\draw[mute,thick] (-0.3,0) -- (9.6,0) node[right,font=\scriptsize,text=sea]{Rx};
% ACK
\draw[arr,teal] (0,0) -- (1.6,1.6);
\node[lab,below] at (0,0) {ACK sent};
\node[lab,above left] at (1.7,1.6) {ACK received:\\$\hat\epsilon$ on preamble,\\outcome of earlier packet};
% processing + wait
\draw[amber,very thick] (1.6,1.6) -- (2.0,1.6);
\draw[sea,very thick,dashed] (2.0,1.6) -- (4.0,1.6);
\node[lab,below] at (3.0,1.6) {wait $h$};
% data packet
\fill[sea!80] (4.0,1.52) rectangle (5.6,1.68);
\node[lab,above] at (4.8,1.6) {data, duration $T_a$};
\fill[teal!25] (4.0,1.6) -- (5.6,1.6) -- (7.2,0) -- (5.6,0) -- cycle;
\fill[sea!80] (5.6,-0.08) rectangle (7.2,0.08);
\node[lab,below] at (6.4,0) {received in\\$[t+h+\tau,\ t+h+\tau+T_a]$};
% horizon brace
\draw[decorate,decoration={brace,mirror,amplitude=4pt}] (2.0,-1.0) -- (7.2,-1.0) node[midway,below=4pt,lab]{forecast beyond the decision time $t$: $h+\tau+T_a$};
\draw[mute,dotted] (2.0,1.6) -- (2.0,-1.0);
\draw[mute,dotted] (7.2,0) -- (7.2,-1.0);
\node[lab,text=amber!85!black,above right] at (1.9,1.6) {$t$: decide $a$};
\end{tikzpicture}

\medskip
\footnotesize\raggedright
$\tau$: one-way propagation time ($\approx$0.67\,s/km). $t_{\mathrm{obs}}$: epoch of the physical state that $\hat\epsilon$ represents, fixed by the measurement equation \open. From an estimate at $t_{\mathrm{obs}}$ the horizon is $H=(t-t_{\mathrm{obs}})+h+\tau+T_a$; if the filter is already propagated to $t$, $h+\tau+T_a$ remains. The ACK of this packet arrives $\ge h+2\tau+T_a$ after the decision, so the margin reacts to old outcomes.
\end{frame}

\begin{frame}{Shorter packets or more frequent tracking?}
\begin{columns}[T]
\begin{column}{0.52\textwidth}
\small
Doppler perfectly corrected at the start of an interval, then held fixed; roughly constant drift, narrowband around $f_c$:
\[
|\Delta f(T)|\approx f_c\,|d|\,T
\quad\Rightarrow\quad
T\le\frac{\Delta f_{\max}}{f_c\,|d|}
\]
$T$ is the interval \alert{without a compensation update}:
\footnotesize
\begin{itemize}
\item blockwise receiver: $T=$ whole packet $\Rightarrow$ shorten $\Nsym$
\item tracking receiver: $T=T_{\mathrm{upd}}=\Kup(T_s+T_{cp})$ $\Rightarrow$ reduce $\Kup$, keep $\Nsym$
\item forecast-aided (Block 1): $|d|\to|d-\hat d|$
\end{itemize}
\end{column}
\begin{column}{0.44\textwidth}
\begin{tcolorbox}[colback=foam!60,colframe=foam!60,arc=2mm,boxrule=0pt]
\footnotesize
\textbf{Illustration only} \sugg\\[1mm]
$f_c=20$\,kHz, assumed $\Delta f_{\max}=2$\,Hz:\\
$|d|=10^{-4}\,$s$^{-1}$ $\Rightarrow$ $T\le1$\,s\\
$|d|=2\times10^{-4}\,$s$^{-1}$ $\Rightarrow$ $T\le0.5$\,s\\[1mm]
$\Delta f_{\max}$ is not a measured receiver threshold.
\end{tcolorbox}
\footnotesize
\textbf{Left out:} wideband resampling effects, differential path Doppler, accumulated phase error, estimator noise, the actual detector. Valid only for fixed correction and roughly constant drift: it explains the mechanism, it is not a reliability guarantee.
\end{column}
\end{columns}
\end{frame}

\begin{frame}{Block 1: receiver pre-warped along the forecast \sugg}
\begin{columns}[T]
\begin{column}{0.52\textwidth}
\small\setlength{\abovedisplayskip}{3pt}\setlength{\belowdisplayskip}{3pt}
Motion time-warps the received signal: $r(t)\approx s\big(t+\int_0^t\epsilon(u)\,du\big)$.

\medskip
\textbf{Today:} estimate $\epsilon$ on the preamble (constant, or constant + acceleration), resample, track.

\smallskip
\textbf{Here:} resample along the \alert{forecast} $\hat\epsilon(t)$ over the whole packet; the tracker sees only the residual warp $\int(\epsilon-\hat\epsilon)$:
\[
|\Delta f(T)|\approx f_c\,|d-\hat d|\,T
\quad\Rightarrow\quad
T\le\frac{\Delta f_{\max}}{f_c\,|d-\hat d|}
\]
{\footnotesize Same simplified model as before, with the drift replaced by the forecast error. Closest: Ceballos Carrascosa--Stojanovic, JOE 2010; Abdelkareem 2016 (gating only).}
\end{column}
\begin{column}{0.44\textwidth}
\footnotesize
\textbf{What it buys}
\begin{itemize}
\item fewer pilots at equal PER, or longer packets, even if acceleration-aware receivers absorb the drift
\item PER now depends on the \alert{forecast uncertainty} $\sigma_e$, which the estimator reports
\end{itemize}
\textbf{What it needs} \open
\begin{itemize}
\item the forecast at node B: a few header bits, or the same filter running there
\item a fallback if the forecast is wrong: residual test, preamble-only
\end{itemize}
\end{column}
\end{columns}
\end{frame}

\begin{frame}{``More pilots'' made explicit}
\begin{columns}[T]
\begin{column}{0.5\textwidth}
\small
Five different things are called pilots:
\begin{enumerate}
\item more frequent \alert{temporal} tracking observations
\item more pilot subcarriers inside one OFDM symbol
\item periodic full training symbols or midambles
\item known preambles at packet boundaries
\item decision-directed observations, if the receiver supports them
\end{enumerate}
\end{column}
\begin{column}{0.46\textwidth}
\footnotesize
Against \emph{temporal} drift the variable that matters is the \alert{spacing between usable tracking observations}, $\Kup$.

\medskip
Baseline: uniform spacing $\Kup\in\{8,4,2\}$, test points, not thresholds. Proposed: a uniform grid at the largest spacing $K_{\max}$ that channel estimation allows, plus extra positions from the forecast uncertainty (Block 2) \sugg. Per-symbol tracking needs embedded pilots and its own overhead model \open.

\medskip
An observation and a receiver update are not the same operation: each receiver must state its estimator and what it needs from an observation \open.
\end{column}
\end{columns}
\end{frame}

\begin{frame}{Block 2: pilots where the forecast is uncertain \sugg}
\begin{columns}[T]
\begin{column}{0.52\textwidth}
\small\setlength{\abovedisplayskip}{3pt}\setlength{\belowdisplayskip}{3pt}
The estimator predicts how the error variance $\sigma_e^2(u)$ grows over the packet, $u\in[0,T_a]$; each training symbol resets it locally.

\smallskip
\textbf{Placement problem}
\[
\min_{\Pset\supseteq\Pset_0:\ |\Pset|\le N_{\mathrm{tr}}}\ \int_0^{T_a}\sigma^2_{e\mid\Pset}(u)\,du
\]
\footnotesize
$\Pset_0$: uniform grid at the largest spacing $K_{\max}$ the tap coherence time allows, so equalisation never loses the channel; the forecast adds only the \alert{extra} pilots. Baseline: $\Pset=\Pset_0$.

\smallskip
\open\ $K_{\max}$ is receiver-dependent (comb pilots relax it; $K_{\max}\ge\Nsym$: $\Pset_0=\emptyset$); integral or max over $u$; not submodular in general (Jawaid--Smith 2015).
\end{column}
\begin{column}{0.44\textwidth}
\footnotesize
\textbf{Expected structure} \claim\ (to derive)\\
For a swell $d=A\sin(\omega t+\varphi)$, a phase error $\delta\varphi$ gives a drift error $A\,\delta\varphi\cos(\omega t+\varphi)$: largest where $d$ crosses zero, where the jerk peaks. Phase error grows with the horizon, so it dominates.

\medskip
$\Rightarrow$ extra pilots cluster around the zero crossings of $d$. The jerk, which hurts acceleration-aware receivers most, decides \alert{where the extra pilots go}.

\smallskip
\textbf{Theory target (R3):} closed-form placement for one swell component; weak-submodularity bound in general; overhead saved vs uniform spacing (see Risks).
\end{column}
\end{columns}
\end{frame}

\begin{frame}{Block 3: both ends measure Doppler \sugg}
\begin{columns}[T]
\begin{column}{0.5\textwidth}
\small\setlength{\abovedisplayskip}{3pt}\setlength{\belowdisplayskip}{3pt}
Node B measures $\hat\epsilon$ on each data packet, node A on each ACK. Both see the same range rate, so a \alert{Kalman filter} on the drift-plus-oscillators model fuses the two interleaved series (an EKF while the sea parameters are learned) \sugg.

\smallskip
Largest gap between samples (stop-and-wait):
\[
\Delta_{\text{one-way}}\ge2\tau+T_a,\qquad \Delta_{\text{two-way}}\approx\tau+T_a
\]
No aliasing for a swell of period $P$ ($\Delta<P/2$):
\[
R<\tfrac c2\big(\tfrac P2-T_a\big)\ \longrightarrow\ R<c\big(\tfrac P2-T_a\big)
\]
{\footnotesize $P=8$\,s, $T_a=1$\,s: 2.25\,km $\to$ 4.5\,km.}
\end{column}
\begin{column}{0.47\textwidth}
\footnotesize
\textbf{Cost and open points} \open
\begin{itemize}\setlength{\itemsep}{0pt}
\item B's estimate returns in the ACK header: value + time stamp, $\approx$20 bits \sugg
\item it arrives with A's own sample and is older: processed in time order; out-of-sequence only with pipelined packets \sugg
\item holds both ends' motion projected on the path, unlike one IMU
\item each direction may lock onto a different path: next slide
\end{itemize}
\textbf{Theory target (R2):} observability and forecast error vs range, one- vs two-way sampling.
\end{column}
\end{columns}
\end{frame}

\begin{frame}{Block 3: what each $\hat\epsilon$ actually measures \sugg}
\begin{columns}[T]
\begin{column}{0.52\textwidth}
\small
Paths have different Doppler scales: surface bounces add wave-induced Doppler, arrival angles differ. The preamble locks onto the first or the strongest arrival, which can \alert{switch} over time and differ between the two directions (beam patterns, depths, surface change over $\tau$).

\medskip
\textbf{Measurement model} for direction $j\in\{A,B\}$:
\[
\hat\epsilon_j(t_k)=\epsilon(t_k)+b_j(t_k)+w_{j,k},\qquad w_{j,k}\sim\mathcal N(0,R_j)
\]
{\footnotesize $\epsilon$: common bulk state; $b_j$: slowly varying bias of the path locked by direction $j$ (a state); $R_j$: per-direction noise, adapted from the innovations.}
\end{column}
\begin{column}{0.44\textwidth}
\footnotesize
\textbf{Path switches}
\begin{itemize}
\item innovation gating: a sample far from the prediction is flagged, not fused blindly
\item on a flagged jump: reset $b_j$ or inflate its covariance, keep the bulk state
\item path-resolved estimates (delay of the locked arrival) as side information, where the receiver gives them
\end{itemize}
\textbf{First check} \open\\
How often $\hat\epsilon$ jumps on Channel Library records, and how far apart are its values on the first and the strongest arrival, before building more machinery.
\end{column}
\end{columns}
\end{frame}

\begin{frame}{Block 3: motion sensors and uplink sampling \sugg}
\begin{columns}[T]
\begin{column}{0.5\textwidth}
\footnotesize
\textbf{Motion sensors, when a node has them}
\begin{itemize}\setlength{\itemsep}{1pt}
\item IMU or GNSS: an extra measurement of the node's \emph{own} motion in the filter
\item the projection onto the propagation path is a state, learned from the link Doppler
\item link Doppler still gives what the sensor cannot: the far end's motion, the path actually used
\item nodes without sensors run on link Doppler alone
\end{itemize}
\textbf{Evaluation:} link only vs sensors only vs both. This answers the obvious question: why not just an IMU?

\smallskip
{\color{mute}Closest: Oiwa et al.\ 2020, accelerometer-aided compensation at one end.}
\end{column}
\begin{column}{0.46\textwidth}
\footnotesize
\textbf{Uplink with $K$ nodes}
\begin{itemize}\setlength{\itemsep}{1pt}
\item the gateway's periodic beacon gives \alert{every node} a Doppler sample at once
\item the gateway measures each node's uplink packet and returns its value in the next beacon ($\approx$20 bits per node)
\item each node runs the two-way filter, with the beacon in place of the ACK
\end{itemize}
\open\ Beacon length grows with $K$; can the nodes share the gateway's own motion?
\end{column}
\end{columns}
\end{frame}

\begin{frame}{Block 4: pre-compensation at the transmitter \sugg}
\begin{columns}[T]
\begin{column}{0.5\textwidth}
\footnotesize\setlength{\abovedisplayskip}{3pt}\setlength{\belowdisplayskip}{3pt}
\textbf{Why only a forecast can do it.} The transmitter's newest $\hat\epsilon$ is at least $\tau$ old; at 2\,km, 2--3\,s from beacon to arrival, about a third of an 8\,s swell period. Using the last value fails; a forecast over the receive interval does not.

\smallskip
\textbf{Mechanism} (first order): node $k$ sends
\[
\tilde s_k(t)=s_k\Big(t-\textstyle\int_0^t\hat\epsilon_k\,du\Big)
\]
\[
\Rightarrow\ r_k(t)\approx s_k\Big(t+\textstyle\int_0^t\big(\epsilon_k-\hat\epsilon_k\big)du\Big)
\]
with $\hat\epsilon_k$ aligned to the receive interval: the gateway sees only the residual.

\smallskip
With $K=1$: an alternative or a complement to Block 1.
\end{column}
\begin{column}{0.46\textwidth}
\footnotesize
\textbf{What it enables: a concurrent uplink}
\begin{itemize}\setlength{\itemsep}{1pt}
\item $K$ nodes in separate subbands, each with its own Doppler
\item at the gateway: one resampling branch per user (Tu et al.\ 2013), \alert{interference remains}
\item pre-compensated at the source: one FFT, per-user residual tracking
\end{itemize}
\textbf{Rough estimate:} users 0.45\,m/s apart at 20\,kHz: $\approx$6\,Hz, half a subcarrier; error 0.05\,m/s: $\approx$0.7\,Hz, interference near $-18$\,dB.

\smallskip
\open\ Only bulk Doppler is removed; path-specific Doppler stays (Zhang 2026).
\end{column}
\end{columns}
\end{frame}

\begin{frame}{Decision: a threshold rule on the forecast, written out \sugg}
\begin{columns}[T]
\begin{column}{0.54\textwidth}
\footnotesize\setlength{\abovedisplayskip}{3pt}\setlength{\belowdisplayskip}{3pt}
\textbf{Predicted PER}: probability, under the forecast distribution of $e$ over the receive interval, that the residual leaves the calibrated tolerance,
\[
\hat p_{\mathrm{fail}}(a,x_t)=\Pr\nolimits_{\hat d,\sigma_e}\big(e\notin\mathcal E_{\mathrm{tol}}(a,\widehat{\mathrm{SNR}})\big)
\]
\[
\mathcal S_t=\big\{a:\ \hat p_{\mathrm{fail}}(a,x_t)\le\alpha-m_t\big\},\quad
a_t=\arg\max_{a\in\mathcal S_t}\tfrac{B_a}{T_a}\big(1-\hat p_{\mathrm{fail}}\big)
\]
\footnotesize
$\alpha$: PER target; $m_t$: margin from delayed ACKs (safety layer). Without the prior ($\hat d=0$): the plain rule predicted drift + margin $\le\dtol$. The PER model is not learned online (follow-up bandit).

\vspace{0.5mm}
The bound must hold over the \emph{whole} receive interval: a pointwise quantile does not bound the maximum. \open\ $\widehat{\mathrm{SNR}}$ source, bits, age (oracle at first).
\end{column}
\begin{column}{0.42\textwidth}
\footnotesize
\textbf{Procedure per packet}
\begin{enumerate}\setlength{\itemsep}{0pt}
\item fuse both ends' Doppler, in time order, with gating
\item forecast $\hat d$, $\sigma_e$ per receive interval
\item pilots per action: grid $\Pset_0$ + extra (Block 2)
\item predicted PER from the calibration curves
\item safe set with $m_t$; if empty: fallback, wait or drop
\item highest expected goodput in the safe set
\item send frame + forecast; B pre-warps
\item update $m_t$ from delayed ACKs
\end{enumerate}
\end{column}
\end{columns}
\end{frame}

\begin{frame}{Uplink: the forecast uncertainty sets guard bands and admission \sugg}
\begin{columns}[T]
\begin{column}{0.52\textwidth}
\footnotesize\setlength{\abovedisplayskip}{3pt}\setlength{\belowdisplayskip}{3pt}
\textbf{Guard band of node $k$}: the offset its residual can reach, at confidence $1-\delta$, over the receive interval,
\[
g_k=\bar f_k\,z_{1-\delta}\,\max_{u}\sigma_{\epsilon,k}(u)
\]
{\scriptsize $\bar f_k$: upper edge of the subband; $\sigma_{\epsilon,k}(u)$: forecast standard deviation of $\epsilon_k$; $z_{1-\delta}$: Gaussian quantile.}

\smallskip
\textbf{Admission and allocation} at the gateway:
\[
\max_{\mathcal K,\{S_k\}}\ \sum_{k\in\mathcal K}\tfrac{B_k}{T}\big(1-\hat p_k\big)
\]
\[
\text{s.t.}\ \sum_{k\in\mathcal K}\big(|S_k|+g_k\big)\le B,\quad \hat p_k\le\alpha-m_{k,t}
\]
Nodes whose guard band costs more than it is worth are deferred to a later slot.
\end{column}
\begin{column}{0.44\textwidth}
\footnotesize
\textbf{What this changes}
\begin{itemize}\setlength{\itemsep}{1pt}
\item $\sigma_e$ gets a measurable effect: how many nodes share the uplink
\item the point-to-point rule is the case $K=1$, $g_1=0$
\item margin $m_{k,t}$ per node, from its own ACKs
\end{itemize}
\textbf{Theory target (R4):} inter-user interference vs residuals and guard bands; nodes supported vs range, sea state, PER target.

\smallskip
{\color{mute}Closest: Hasan et al.\ 2025 (no forecast).} \open\ Path-specific Doppler.
\end{column}
\end{columns}
\end{frame}

% ===========================================================
\section{Theoretical objectives}

\begin{frame}{Theory: four core results, the rest as extensions or follow-up}
\small
\begin{columns}[T,onlytextwidth]
\begin{column}{0.24\textwidth}
\begin{block}{R1 Predictability}
\scriptsize
How long Doppler stays predictable (L1; idea document: $\approx$13\,s for an 8\,s swell vs $\approx$1\,s for fading); mixed seas degrade gracefully (P2)
\end{block}
\end{column}
\begin{column}{0.24\textwidth}
\begin{block}{R2 Sampling}
\scriptsize
Forecast error vs range, one- vs two-way ($\times$2 samples per round trip); observability with delayed samples and per-direction bias (L3--L4); gain from motion sensors
\end{block}
\end{column}
\begin{column}{0.24\textwidth}
\begin{block}{R3 Pilots}
\scriptsize
Optimal extra positions on top of the grid $\Pset_0$ for one swell component; greedy bound only via weak submodularity (from P7)
\end{block}
\end{column}
\begin{column}{0.24\textwidth}
\begin{block}{R4 Concurrency}
\scriptsize
Inter-user interference vs forecast residuals and guard bands; how many nodes share the uplink at a given range, sea state and PER target
\end{block}
\end{column}
\end{columns}

\smallskip
\footnotesize
\textbf{Weight:} R1 and R2 carry the paper, R4 is the main application; R3 if the single-swell case closes. \textbf{Follow-up:} bandit regret with arm-dependent delay (F1).\\
\textbf{Supporting:} T10--T11, PER safety layer and its stability (known structure: OLLA / ACI).\\
\textbf{Extensions:} T8 quiet windows, T9 Doppler vs SNR feedback, L3--L4 path separation, P5, P6.

\smallskip
\open\ Hard parts: path switches and bias in R2; submodularity in R3; path-specific Doppler in R4. {\color{mute}L/P/T: results of the idea document.}
\end{frame}

\begin{frame}{Safety layer: what makes the PER claim non-vacuous}
\begin{columns}[T]
\begin{column}{0.5\textwidth}
\footnotesize
\textbf{Deck claim} \claim: long-run PER below target for any channel, predictor and estimate.

\medskip
\textbf{Margin update.} The deck's verbal rule, written as
\[
m_{k+1}=m_k+\eta\,(e_k-\alpha),\quad e_k=\mathbf 1\{\text{failure}\}
\]
per processed outcome ($\alpha$: PER target). Unprojected, $\frac1N\sum_k e_k=\alpha+\frac{m_{N+1}-m_1}{\eta N}$: calibration only if $m$ grows sublinearly (if all fail, it does not). \open\ Signed correction or buffer $m\ge0$? (projection changes the proof); outcome vs packet index.

\medskip
\textbf{Which guarantee?} Long-run frequency, finite-horizon, or per-packet. Report it with goodput, drops, deadline misses, retransmissions.
\end{column}
\begin{column}{0.46\textwidth}
\footnotesize
\textbf{To fix before proving Thm.\ 10} \open
\begin{itemize}
\item PER denominator: packets sent? with retransmissions? deadline expiries?
\item missing ACK or timeout: failure, or excluded?
\item empty feasible set: fallback action, wait, or drop, and how each is counted
\item minimum service required, e.g.\ a bound on the dropped or delayed fraction
\item feasibility assumption, e.g.\ some action always meets the target with slack
\item action-dependent delay: outcomes arrive $\ge h+2\tau+T_a$ late
\end{itemize}
\end{column}
\end{columns}
\end{frame}

\begin{frame}{Risks found by the prior-art search}
\begin{columns}[T]
\begin{column}{0.5\textwidth}
\footnotesize
\textbf{R3: no free greedy bound.} Kalman sensor-scheduling costs are not submodular in general (Jawaid--Smith, Automatica 2015); the Krause et al.\ 2008 bound covers Gaussian-process placement only. Routes: closed form for one swell; submodularity ratio or curvature (Hashemi et al., TAC 2020); Jawaid--Smith conditions for our model.

\medskip
\textbf{Block 2: periodic is optimal when stationary.} Uniform pilots minimise steady-state MSE for a Gauss--Markov channel (Dong--Tong--Sadler 2002). Extras pay off only where $\sigma_e^2(u)$ varies within the packet: state it as a precondition.
\end{column}
\begin{column}{0.46\textwidth}
\footnotesize
\textbf{Block 1: field counter-evidence.} Surface-buoy trial, 10.5\,h, 26\,170 frames: per-symbol Doppler refinement lowers PER by at most $1.15\times$ (Lee, Park, Kim, Im, SSRN 2026). Block 1 must be shown in a regime fixed in advance (long packets, strong swell, few pilots), against acceleration-aware baselines (Zakharov 2019, Cho--Kim 2026, Hu 2024).

\medskip
\textbf{Framing.} DE-Sync's exchange, navigation filters on modem Doppler and NTN make ``close the loop'' look obvious. Defence: the analysis (R1 horizon, R2 two-way sampling with bias, R3 placement, R4 concurrency), not the plumbing.
\end{column}
\end{columns}
\end{frame}

% ===========================================================
\section{Follow-up}

\begin{frame}{Follow-up paper: a contextual bandit that knows the physics \sugg}
\begin{columns}[T]
\begin{column}{0.52\textwidth}
\footnotesize\setlength{\abovedisplayskip}{3pt}\setlength{\belowdisplayskip}{3pt}
{\color{mute}Not in the JSAC submission: it replaces the threshold rule once the calibration curves are learned online.}

\smallskip
\textbf{Context} $x_t$: forecast over the candidate receive interval ($\hat d$, $\sigma_e$), $\widehat{\mathrm{SNR}}$, range.

\smallskip
\textbf{Arms} $a=(M,\Nsym)$; for each arm, $\Pset$ comes from Block 2.

\smallskip
\textbf{Reward}: delivered bits per unit airtime, $B_a\,\mathbf 1\{\text{success}\}/T_a$ (one objective, fixed before the comparison).

\smallskip
\textbf{Model}: failure probability monotone in forecast error and action; generalised linear in physical features,
\[
p_{\mathrm{fail}}(a,x)=g\big(\theta_a^\top\psi(x)\big),
\]
$\theta_a$ initialised from the calibration curves, learned online (UCB or Thompson sampling).
\end{column}
\begin{column}{0.44\textwidth}
\footnotesize
\textbf{What is new for the theory (F1)}
\begin{itemize}
\item feedback for arm $a$ arrives $D_a\ge h+2\tau+T_a$ later: the delay \alert{depends on the arm}
\item short packets return feedback sooner, biasing a naive learner towards them \open
\item target \claim: regret $\tilde O(\sqrt{T})$ + delay term, extending fixed-delay results (Joulani et al., ICML 2013)
\end{itemize}
\textbf{Safety layer}: only arms meeting the PER target with margin $m_t$.

{\color{mute} Builds on AMUSE, POSEIDON, AoI bandit (Busacca et al.). Details: \texttt{followup\_bandit.md}.}
\end{column}
\end{columns}
\end{frame}

% ===========================================================
\section{Questions for the group}

\begin{frame}{Questions for the group}
\begin{columns}[T]
\begin{column}{0.5\textwidth}
\small
\begin{itemize}
\item Bandit moved to a follow-up: is a calibrated threshold rule enough for the JSAC paper?
\item Which receiver structure sets $K_{\max}$: comb pilots in every symbol, or full training symbols only?
\item Can the target modem or simulator carry $\approx$20 header bits for the forecast and the far end's Doppler, and $K$ such fields in the beacon?
\item What is already implemented, and which results are proved?
\end{itemize}
\end{column}
\begin{column}{0.46\textwidth}
\small
\begin{itemize}
\item May the policy abstain, drop packets or use a fallback, and how is service counted?
\item How much time and theory/simulation support is available?
\item Target uplink: buoy gateway with seabed nodes, AUVs, or both? Which nodes have an IMU or GNSS?
\item First test: does a forecast prior still help acceleration-aware receivers, in PER or in pilots saved?
\end{itemize}
\end{column}
\end{columns}
\end{frame}

\end{document}
